\documentclass[tikz, border=10mm]{standalone}
\usepackage{tikz}
\usetikzlibrary{positioning, calc, arrows.meta, chains, fit, shapes, 3d}
% nn_blocks.tex
% Defines styles for neural network blocks using TikZ
% Now supports Hybrid 2D/3D Visualization with PlotNeuralNet

% --- PlotNeuralNet Integration ---
% We assume Ball.sty, Box.sty, RightBandedBox.sty are in the same directory
\usepackage{Ball}
\usepackage{Box}
\usepackage{RightBandedBox}

% PlotNeuralNet Init Definitions (from init.tex)
\usetikzlibrary{quotes,arrows.meta}
\usetikzlibrary{positioning}
\def\edgecolor{rgb:blue,4;red,1;green,4;black,3}
\newcommand{\midarrow}{\tikz \draw[-Stealth,line width =0.8mm,draw=\edgecolor] (-0.3,0) -- ++(0.3,0);}

% Define PlotNeuralNet Colors (Standard Palette)
\definecolor{ConvColor}{rgb}{0.98,0.81,0.69}  % Apricot/Orange
\definecolor{ConvReluColor}{rgb}{0.98,0.81,0.69}
\definecolor{PoolColor}{rgb}{0.8,0.33,0.33}   % Redish
\definecolor{UnpoolColor}{rgb}{0.6,0.7,0.9}   % Blueish
\definecolor{FcColor}{rgb}{0.6,0.2,0.8}       % Purple
\definecolor{FcReluColor}{rgb}{0.6,0.2,0.8}
\definecolor{SoftmaxColor}{rgb}{0.9,0.9,0.6}  % Yellowish
\definecolor{SumColor}{rgb}{0.9,0.9,0.9}      % Gray

% Custom colors for our specific blocks
\definecolor{ResBlockColor}{rgb}{0.85,0.92,1.0} % Light Blue
\definecolor{EmbedColor}{rgb}{0.95,0.95,0.8}    % Light Yellow
\definecolor{TransColor}{rgb}{0.8,0.9,0.8}      % Light Green

% --- NNTikZ 2D Styles (Matched to EDSR Style) ---
\usetikzlibrary{positioning, chains, shapes.geometric, fit, shapes, arrows.meta, calc, backgrounds, shadows}

\tikzset{
    >=LaTeX, 
    very thick,
    node distance=1.0cm and 1.2cm,
    % Basic Arrow
    arrow/.style={
        ->,
        very thick,
        draw=black!80,
        rounded corners=0.2cm
    },
    % Dashed Skip Connection
    skip/.style={
        arrow,
        dashed,
        draw=blue!60
    },
    % Grouping Block (Background Box)
    block/.style={
        rectangle,
        fill=gray!10,
        rounded corners=3mm,
        draw=black!60,
        very thick,
        inner sep=0.8em,
        drop shadow={opacity=0.15}
    },
    nntikz_block/.style={block}, % Alias
    % General Layer
    layer/.style={
        rectangle,
        fill=white!10,
        rounded corners=1mm,
        inner xsep=0.6em,
        inner ysep=0.6em,
        minimum height=3.0em,
        align=center,
        draw=black!80,
        very thick,
        drop shadow={opacity=0.2, shadow xshift=1mm, shadow yshift=-1mm}
    },
    nntikz_layer/.style={layer}, % Alias
    % Specific Layer Types
    conv/.style={
        layer,
        fill=orange!10,
        draw=orange!80!black
    },
    relu/.style={
        layer,
        fill=yellow!10,
        draw=yellow!80!black,
        minimum height=2.5em
    },
    pool/.style={
        layer,
        fill=red!10,
        draw=red!80!black
    },
    up/.style={
        layer,
        fill=purple!10,
        draw=purple!80!black
    },
    res/.style={ % Residual Block or similar
        layer,
        fill=blue!5,
        draw=blue!80!black
    },
    op/.style={ % Other operations
        layer,
        fill=green!10,
        draw=green!80!black
    },
    % Input/Output
    input/.style={ 
        circle,
        minimum width=3.0em,
        draw=black!80,
        fill=gray!20,
        thick,
        align=center,
        drop shadow={opacity=0.2}
    },
    io/.style={input}, % Alias
    % Summation
    sum/.style={
        circle,
        draw=black!80,
        fill=white,
        inner sep=0pt,
        minimum size=1.2em,
        node contents={+}
    }
}



\definecolor{convcolor}{rgb}{1,0.8,0.5}
\definecolor{rescolor}{rgb}{0.8,0.8,1}
\definecolor{poolcolor}{rgb}{1,0.5,0.5}
\definecolor{upcolor}{rgb}{0.8,0.5,1}
\definecolor{opcolor}{rgb}{0.5,1,0.5}

\def\ConvColor{convcolor}
\def\ResColor{rescolor}
\def\PoolColor{poolcolor}
\def\UpColor{upcolor}
\def\OpColor{opcolor}
\def\ConvColor{convcolor}
\def\ResColor{rescolor}
\def\OpColor{opcolor}
\def\PoolColor{poolcolor}

\begin{document}
\begin{tikzpicture}

% 1. Input Image (Branch Input)
\node[io] (in) {Input};

% 2. Branch Net (CNN Backbone)
\pic[shift={(2.0 * 4.300000000000001,0,0)}] at (in.east) {RightBandedBox={
    name=b_conv1,
    caption=Branch\\Conv1,
    xlabel={"64", ""},
    zlabel=64,
    fill=\ConvColor,
    height=30, width=2, depth=30
}};

\pic[shift={(1.5 * 4.300000000000001,0,0)}] at (b_conv1-east) {RightBandedBox={
    name=b_conv2,
    caption=Branch\\Conv2,
    xlabel={"128", ""},
    zlabel=32,
    fill=\ConvColor,
    height=20, width=3, depth=20
}};

\pic[shift={(1.5 * 4.300000000000001,0,0)}] at (b_conv2-east) {RightBandedBox={
    name=b_conv3,
    caption=Branch\\Conv3,
    xlabel={"256", ""},
    zlabel=16,
    fill=\ConvColor,
    height=15, width=4, depth=15
}};

% Global Pooling + FC -> Coeffs
\pic[shift={(1.5 * 4.300000000000001,0,0)}] at (b_conv3-east) {Box={
    name=b_gap,
    caption=GAP,
    fill=\PoolColor,
    height=5, width=1, depth=5
}};

\pic[shift={(1.2 * 4.300000000000001,0,0)}] at (b_gap-east) {Box={
    name=b_vec,
    caption=Coeffs\\$b$,
    fill=\ResColor,
    height=1, width=6, depth=1
}};

% 3. Trunk Net (MLP)
% Coords input (below Branch Conv1)
\node[io, below=12.0cm of b_conv1-west] (coords) {Coords};

\pic[shift={(2.0 * 4.300000000000001,0,0)}] at (coords.east) {Box={
    name=t_mlp1,
    caption=Trunk\\MLP1,
    fill=\ConvColor,
    height=8, width=2, depth=8
}};

\pic[shift={(1.0 * 4.300000000000001,0,0)}] at (t_mlp1-east) {Box={
    name=t_mlp2,
    caption=Trunk\\MLP2,
    fill=\ConvColor,
    height=8, width=2, depth=8
}};

\pic[shift={(1.0 * 4.300000000000001,0,0)}] at (t_mlp2-east) {Box={
    name=t_mlp3,
    caption=Trunk\\MLP3,
    fill=\ConvColor,
    height=8, width=2, depth=8
}};

% Basis
\pic[shift={(1.5 * 4.300000000000001,0,0)}] at (t_mlp3-east) {Box={
    name=t_vec,
    caption=Basis\\$t$,
    fill=\ResColor,
    height=1, width=6, depth=1
}};

% 4. Fusion (Dot Product)
% Position fusion roughly between the ends of both branches
\coordinate (fusion_pos) at ($(b_vec-east)!0.5!(t_vec-east) + (6.0, 0, 0)$);

\pic[shift={(0,0,0)}] at (fusion_pos) {Box={
    name=dot,
    caption=Dot,
    fill=\OpColor,
    height=5, width=2, depth=5
}};

% 5. Output
\node[io, right=1.5cm of dot-east] (out) {Output};

% Connections
% Branch
\draw[arrow] (in) -- (b_conv1-west);
\draw[arrow] (b_conv1-east) -- (b_conv2-west);
\draw[arrow] (b_conv2-east) -- (b_conv3-west);
\draw[arrow] (b_conv3-east) -- (b_gap-west);
\draw[arrow] (b_gap-east) -- (b_vec-west);

% Trunk
\draw[arrow] (coords) -- (t_mlp1-west);
\draw[arrow] (t_mlp1-east) -- (t_mlp2-west);
\draw[arrow] (t_mlp2-east) -- (t_mlp3-west);
\draw[arrow] (t_mlp3-east) -- (t_vec-west);

% Fusion
\draw[arrow] (b_vec-east) -- ++(1.5,0,0) |- (dot-west);
\draw[arrow] (t_vec-east) -- ++(1.5,0,0) |- (dot-west);

% Output
\draw[arrow] (dot-east) -- (out);

\end{tikzpicture}
\end{document}